% Options for packages loaded elsewhere
% Options for packages loaded elsewhere
\PassOptionsToPackage{unicode}{hyperref}
\PassOptionsToPackage{hyphens}{url}
\PassOptionsToPackage{dvipsnames,svgnames,x11names}{xcolor}
%
\documentclass[
  11pt,
  a4paper,
  DIV=11,
  numbers=noendperiod]{scrreprt}
\usepackage{xcolor}
\usepackage[top=24mm,bottom=24mm,left=22mm,right=22mm]{geometry}
\usepackage{amsmath,amssymb}
\setcounter{secnumdepth}{5}
\usepackage{iftex}
\ifPDFTeX
  \usepackage[T1]{fontenc}
  \usepackage[utf8]{inputenc}
  \usepackage{textcomp} % provide euro and other symbols
\else % if luatex or xetex
  \usepackage{unicode-math} % this also loads fontspec
  \defaultfontfeatures{Scale=MatchLowercase}
  \defaultfontfeatures[\rmfamily]{Ligatures=TeX,Scale=1}
\fi
\usepackage{lmodern}
\ifPDFTeX\else
  % xetex/luatex font selection
\fi
% Use upquote if available, for straight quotes in verbatim environments
\IfFileExists{upquote.sty}{\usepackage{upquote}}{}
\IfFileExists{microtype.sty}{% use microtype if available
  \usepackage[]{microtype}
  \UseMicrotypeSet[protrusion]{basicmath} % disable protrusion for tt fonts
}{}
\makeatletter
\@ifundefined{KOMAClassName}{% if non-KOMA class
  \IfFileExists{parskip.sty}{%
    \usepackage{parskip}
  }{% else
    \setlength{\parindent}{0pt}
    \setlength{\parskip}{6pt plus 2pt minus 1pt}}
}{% if KOMA class
  \KOMAoptions{parskip=half}}
\makeatother
% Make \paragraph and \subparagraph free-standing
\makeatletter
\ifx\paragraph\undefined\else
  \let\oldparagraph\paragraph
  \renewcommand{\paragraph}{
    \@ifstar
      \xxxParagraphStar
      \xxxParagraphNoStar
  }
  \newcommand{\xxxParagraphStar}[1]{\oldparagraph*{#1}\mbox{}}
  \newcommand{\xxxParagraphNoStar}[1]{\oldparagraph{#1}\mbox{}}
\fi
\ifx\subparagraph\undefined\else
  \let\oldsubparagraph\subparagraph
  \renewcommand{\subparagraph}{
    \@ifstar
      \xxxSubParagraphStar
      \xxxSubParagraphNoStar
  }
  \newcommand{\xxxSubParagraphStar}[1]{\oldsubparagraph*{#1}\mbox{}}
  \newcommand{\xxxSubParagraphNoStar}[1]{\oldsubparagraph{#1}\mbox{}}
\fi
\makeatother

\usepackage{color}
\usepackage{fancyvrb}
\newcommand{\VerbBar}{|}
\newcommand{\VERB}{\Verb[commandchars=\\\{\}]}
\DefineVerbatimEnvironment{Highlighting}{Verbatim}{commandchars=\\\{\}}
% Add ',fontsize=\small' for more characters per line
\newenvironment{Shaded}{}{}
\newcommand{\AlertTok}[1]{\textcolor[rgb]{1.00,0.33,0.33}{\textbf{#1}}}
\newcommand{\AnnotationTok}[1]{\textcolor[rgb]{0.42,0.45,0.49}{#1}}
\newcommand{\AttributeTok}[1]{\textcolor[rgb]{0.84,0.23,0.29}{#1}}
\newcommand{\BaseNTok}[1]{\textcolor[rgb]{0.00,0.36,0.77}{#1}}
\newcommand{\BuiltInTok}[1]{\textcolor[rgb]{0.84,0.23,0.29}{#1}}
\newcommand{\CharTok}[1]{\textcolor[rgb]{0.01,0.18,0.38}{#1}}
\newcommand{\CommentTok}[1]{\textcolor[rgb]{0.42,0.45,0.49}{#1}}
\newcommand{\CommentVarTok}[1]{\textcolor[rgb]{0.42,0.45,0.49}{#1}}
\newcommand{\ConstantTok}[1]{\textcolor[rgb]{0.00,0.36,0.77}{#1}}
\newcommand{\ControlFlowTok}[1]{\textcolor[rgb]{0.84,0.23,0.29}{#1}}
\newcommand{\DataTypeTok}[1]{\textcolor[rgb]{0.84,0.23,0.29}{#1}}
\newcommand{\DecValTok}[1]{\textcolor[rgb]{0.00,0.36,0.77}{#1}}
\newcommand{\DocumentationTok}[1]{\textcolor[rgb]{0.42,0.45,0.49}{#1}}
\newcommand{\ErrorTok}[1]{\textcolor[rgb]{1.00,0.33,0.33}{\underline{#1}}}
\newcommand{\ExtensionTok}[1]{\textcolor[rgb]{0.84,0.23,0.29}{\textbf{#1}}}
\newcommand{\FloatTok}[1]{\textcolor[rgb]{0.00,0.36,0.77}{#1}}
\newcommand{\FunctionTok}[1]{\textcolor[rgb]{0.44,0.26,0.76}{#1}}
\newcommand{\ImportTok}[1]{\textcolor[rgb]{0.01,0.18,0.38}{#1}}
\newcommand{\InformationTok}[1]{\textcolor[rgb]{0.42,0.45,0.49}{#1}}
\newcommand{\KeywordTok}[1]{\textcolor[rgb]{0.84,0.23,0.29}{#1}}
\newcommand{\NormalTok}[1]{\textcolor[rgb]{0.14,0.16,0.18}{#1}}
\newcommand{\OperatorTok}[1]{\textcolor[rgb]{0.14,0.16,0.18}{#1}}
\newcommand{\OtherTok}[1]{\textcolor[rgb]{0.44,0.26,0.76}{#1}}
\newcommand{\PreprocessorTok}[1]{\textcolor[rgb]{0.84,0.23,0.29}{#1}}
\newcommand{\RegionMarkerTok}[1]{\textcolor[rgb]{0.42,0.45,0.49}{#1}}
\newcommand{\SpecialCharTok}[1]{\textcolor[rgb]{0.00,0.36,0.77}{#1}}
\newcommand{\SpecialStringTok}[1]{\textcolor[rgb]{0.01,0.18,0.38}{#1}}
\newcommand{\StringTok}[1]{\textcolor[rgb]{0.01,0.18,0.38}{#1}}
\newcommand{\VariableTok}[1]{\textcolor[rgb]{0.89,0.38,0.04}{#1}}
\newcommand{\VerbatimStringTok}[1]{\textcolor[rgb]{0.01,0.18,0.38}{#1}}
\newcommand{\WarningTok}[1]{\textcolor[rgb]{1.00,0.33,0.33}{#1}}

\usepackage{longtable,booktabs,array}
\newcounter{none} % for unnumbered tables
\usepackage{calc} % for calculating minipage widths
% Correct order of tables after \paragraph or \subparagraph
\usepackage{etoolbox}
\makeatletter
\patchcmd\longtable{\par}{\if@noskipsec\mbox{}\fi\par}{}{}
\makeatother
% Allow footnotes in longtable head/foot
\IfFileExists{footnotehyper.sty}{\usepackage{footnotehyper}}{\usepackage{footnote}}
\makesavenoteenv{longtable}
\usepackage{graphicx}
\makeatletter
\newsavebox\pandoc@box
\newcommand*\pandocbounded[1]{% scales image to fit in text height/width
  \sbox\pandoc@box{#1}%
  \Gscale@div\@tempa{\textheight}{\dimexpr\ht\pandoc@box+\dp\pandoc@box\relax}%
  \Gscale@div\@tempb{\linewidth}{\wd\pandoc@box}%
  \ifdim\@tempb\p@<\@tempa\p@\let\@tempa\@tempb\fi% select the smaller of both
  \ifdim\@tempa\p@<\p@\scalebox{\@tempa}{\usebox\pandoc@box}%
  \else\usebox{\pandoc@box}%
  \fi%
}
% Set default figure placement to htbp
\def\fps@figure{htbp}
\makeatother





\setlength{\emergencystretch}{3em} % prevent overfull lines

\providecommand{\tightlist}{%
  \setlength{\itemsep}{0pt}\setlength{\parskip}{0pt}}



 


\usepackage{fontspec}
\usepackage{unicode-math}
\usepackage{fancyhdr}
\usepackage{xcolor}
\usepackage{float}
\usepackage{booktabs}
\usepackage{longtable}
\usepackage{array}
\usepackage{tcolorbox}
\usepackage{fvextra}
\definecolor{ForecastBlue}{HTML}{1F4E79}
\definecolor{ForecastTeal}{HTML}{167D8D}
\definecolor{ForecastGold}{HTML}{D99A2B}
\definecolor{ForecastLight}{HTML}{EEF5F8}
\setmainfont{Noto Sans}
\setsansfont{Noto Sans}
\setmonofont{Noto Sans Mono}
\setmathfont{Latin Modern Math}
\pagestyle{fancy}
\fancyhf{}
\fancyfoot[C]{\footnotesize\color{ForecastBlue} Business Forecasting Notes by Dr. Sankalp Naik}
\fancyfoot[R]{\footnotesize\thepage}
\renewcommand{\headrulewidth}{0pt}
\renewcommand{\footrulewidth}{0.4pt}
\addtokomafont{chapter}{\color{ForecastBlue}}
\addtokomafont{section}{\color{ForecastTeal}}
\addtokomafont{subsection}{\color{ForecastBlue}}
\AtBeginDocument{\renewcommand{\contentsname}{Index / Contents}}
\KOMAoption{captions}{tableheading}
\makeatletter
\@ifpackageloaded{tcolorbox}{}{\usepackage[skins,breakable]{tcolorbox}}
\@ifpackageloaded{fontawesome5}{}{\usepackage{fontawesome5}}
\definecolor{quarto-callout-color}{HTML}{909090}
\definecolor{quarto-callout-note-color}{HTML}{0758E5}
\definecolor{quarto-callout-important-color}{HTML}{CC1914}
\definecolor{quarto-callout-warning-color}{HTML}{EB9113}
\definecolor{quarto-callout-tip-color}{HTML}{00A047}
\definecolor{quarto-callout-caution-color}{HTML}{FC5300}
\definecolor{quarto-callout-color-frame}{HTML}{acacac}
\definecolor{quarto-callout-note-color-frame}{HTML}{4582ec}
\definecolor{quarto-callout-important-color-frame}{HTML}{d9534f}
\definecolor{quarto-callout-warning-color-frame}{HTML}{f0ad4e}
\definecolor{quarto-callout-tip-color-frame}{HTML}{02b875}
\definecolor{quarto-callout-caution-color-frame}{HTML}{fd7e14}
\makeatother
\makeatletter
\@ifpackageloaded{caption}{}{\usepackage{caption}}
\AtBeginDocument{%
\ifdefined\contentsname
  \renewcommand*\contentsname{Table of contents}
\else
  \newcommand\contentsname{Table of contents}
\fi
\ifdefined\listfigurename
  \renewcommand*\listfigurename{List of Figures}
\else
  \newcommand\listfigurename{List of Figures}
\fi
\ifdefined\listtablename
  \renewcommand*\listtablename{List of Tables}
\else
  \newcommand\listtablename{List of Tables}
\fi
\ifdefined\figurename
  \renewcommand*\figurename{Figure}
\else
  \newcommand\figurename{Figure}
\fi
\ifdefined\tablename
  \renewcommand*\tablename{Table}
\else
  \newcommand\tablename{Table}
\fi
}
\@ifpackageloaded{float}{}{\usepackage{float}}
\floatstyle{ruled}
\@ifundefined{c@chapter}{\newfloat{codelisting}{h}{lop}}{\newfloat{codelisting}{h}{lop}[chapter]}
\floatname{codelisting}{Program}
\newcommand*\listoflistings{\listof{codelisting}{List of Listings}}
\makeatother
\makeatletter
\makeatother
\makeatletter
\@ifpackageloaded{caption}{}{\usepackage{caption}}
\@ifpackageloaded{subcaption}{}{\usepackage{subcaption}}
\makeatother
\usepackage{bookmark}
\IfFileExists{xurl.sty}{\usepackage{xurl}}{} % add URL line breaks if available
\urlstyle{same}
\hypersetup{
  pdftitle={Classical Forecasting Techniques},
  pdfauthor={Dr.~Sankalp Purushottam Naik},
  colorlinks=true,
  linkcolor={RoyalBlue},
  filecolor={Maroon},
  citecolor={RoyalBlue},
  urlcolor={RoyalBlue},
  pdfcreator={LaTeX via pandoc}}


\title{Classical Forecasting Techniques}
\usepackage{etoolbox}
\makeatletter
\providecommand{\subtitle}[1]{% add subtitle to \maketitle
  \apptocmd{\@title}{\par {\large #1 \par}}{}{}
}
\makeatother
\subtitle{Business Forecasting Notes}
\author{Dr.~Sankalp Purushottam Naik}
\date{2026-09-28}
\begin{document}
\maketitle

\thispagestyle{empty}
\begin{center}
\vspace*{2.8cm}
{\Huge\bfseries\color{ForecastBlue} Classical Forecasting Techniques\par}
\vspace{0.7cm}
{\Large Business Forecasting Notes\par}
\vspace{1.8cm}
{\Large\bfseries Dr. Sankalp Purushottam Naik\par}
\vfill
{\large Data Generating Processes • Naïve Methods • Moving Averages • Holt • Holt-Winters\par}
\vspace{1.0cm}
{\small Prepared as instructional notes for Business Forecasting\par}
\end{center}
\clearpage

\renewcommand*\contentsname{Contents}
{
\setcounter{tocdepth}{2}
\tableofcontents
}

\begin{tcolorbox}[enhanced jigsaw, arc=.35mm, bottomrule=.15mm, bottomtitle=1mm, breakable, colback=white, colbacktitle=quarto-callout-note-color!10!white, colframe=quarto-callout-note-color-frame, coltitle=black, left=2mm, leftrule=.75mm, opacityback=0, opacitybacktitle=0.6, rightrule=.15mm, title=\textcolor{quarto-callout-note-color}{\faInfo}\hspace{0.5em}{📚 Source Attribution \& Disclaimer}, titlerule=0mm, toprule=.15mm, toptitle=1mm]

These notes are prepared for \textbf{teaching and classroom use}. The
conceptual treatment of forecasting methods draws substantially on
\emph{Forecasting: Principles and Practice} by \textbf{Rob J Hyndman and
George Athanasopoulos}. The explanations, examples, Python
implementations, classroom analogies, tables, and presentation structure
in these notes have been adapted and developed for instructional
purposes. Students are encouraged to consult the original text for a
comprehensive treatment of forecasting theory and methods.

\end{tcolorbox}

Classical forecasting techniques are a set of methods used to predict
future values based on historical data. These techniques are widely used
in various fields, including finance, economics, and supply chain
management. In this section, we will explore some of the most commonly
used classical forecasting techniques, including naive methods, moving
average methods, exponential smoothing methods, Holt method, and
Holt-Winters methods. We will also discuss their advantages and
limitations, as well as how to implement them in practice.

\chapter{Data Generating Process
(DGP)}\label{data-generating-process-dgp}

A data generating process (DGP) is a statistical model that describes
how data is generated. It is a theoretical construct that helps us
understand the underlying patterns and relationships in the data. In
practice, we often do not know the true DGP, and we must rely on
assumptions and approximations to make predictions.

Examples of DGPs include linear regression models, time series models,
and stochastic processes. Understanding the DGP is crucial for selecting
appropriate forecasting techniques and for interpreting the results of
forecasts.

\section{True DGP vs.~Models}\label{true-dgp-vs.-models}

True DGP refers to the actual process that generates the data, while
models are simplified representations of the DGP that we use for
forecasting. Models may not capture all the complexities of the true
DGP, but they can still provide useful insights and predictions. It is
important to evaluate the performance of models against the true DGP to
ensure their accuracy and reliability.

\section{Illustration 1: DGP and Models - Study Hours and Exam
Scores}\label{illustration-1-dgp-and-models---study-hours-and-exam-scores}

Suppose we observe the following data for 10 students. Each student
reports the number of hours they studied before an examination, and we
observe their actual examination score.

Student \textbar{} Study Hours \textbar{} Actual Score \textbar{}

------: \textbar{} ----------: \textbar{} -----------: \textbar{}

~~~~~~1 \textbar{} 1.5 \textbar{} 43 \textbar{}

~~~~~~2 \textbar{} 2.0 \textbar{} 47 \textbar{}

~~~~~~3 \textbar{} 3.5 \textbar{} 56 \textbar{}

~~~~~~4 \textbar{} 4.0 \textbar{} 59 \textbar{}

~~~~~~5 \textbar{} 5.5 \textbar{} 68 \textbar{}

~~~~~~6 \textbar{} 6.0 \textbar{} 70 \textbar{}

~~~~~~7 \textbar{} 7.5 \textbar{} 78 \textbar{}

~~~~~~8 \textbar{} 8.0 \textbar{} 81 \textbar{}

~~~~~~9 \textbar{} 9.5 \textbar{} 87 \textbar{}

~~~~~10 \textbar{} 10.0 \textbar{} 89 \textbar{}

At this stage, ****we make no assumptions about how the data were
generated****.

We simply observe the relationship between study hours and actual
scores.

\subsection{\texorpdfstring{\textbf{Data Display and Scatter
Plot}}{Data Display and Scatter Plot}}\label{data-display-and-scatter-plot}

\textbf{::: \{.callout-tip title=``💻 Coding Session''\}}

\begin{Shaded}
\begin{Highlighting}[]
\CommentTok{\#| output: true}

\CommentTok{\#| label: lst{-}dgp{-}1}

\CommentTok{\#| lst{-}cap: "Data Display and Scatter Plot"}

\ImportTok{import}\NormalTok{ pandas }\ImportTok{as}\NormalTok{ pd}

\ImportTok{import}\NormalTok{ matplotlib.pyplot }\ImportTok{as}\NormalTok{ plt}

\CommentTok{\# Create the dataset}

\NormalTok{df }\OperatorTok{=}\NormalTok{ pd.DataFrame(}
\NormalTok{    \{}
        \StringTok{"Student"}\NormalTok{: }\BuiltInTok{range}\NormalTok{(}\DecValTok{1}\NormalTok{, }\DecValTok{11}\NormalTok{),}
        \StringTok{"Study\_Hours"}\NormalTok{: [}\FloatTok{1.5}\NormalTok{, }\FloatTok{2.0}\NormalTok{, }\FloatTok{3.5}\NormalTok{, }\FloatTok{4.0}\NormalTok{, }\FloatTok{5.5}\NormalTok{, }\FloatTok{6.0}\NormalTok{, }\FloatTok{7.5}\NormalTok{, }\FloatTok{8.0}\NormalTok{, }\FloatTok{9.5}\NormalTok{, }\FloatTok{10.0}\NormalTok{],}
        \StringTok{"Actual\_Score"}\NormalTok{: [}\DecValTok{43}\NormalTok{, }\DecValTok{47}\NormalTok{, }\DecValTok{56}\NormalTok{, }\DecValTok{59}\NormalTok{, }\DecValTok{68}\NormalTok{, }\DecValTok{70}\NormalTok{, }\DecValTok{78}\NormalTok{, }\DecValTok{81}\NormalTok{, }\DecValTok{87}\NormalTok{, }\DecValTok{89}\NormalTok{],}
\NormalTok{    \}}
\NormalTok{)}

\CommentTok{\# Display the data}

\BuiltInTok{print}\NormalTok{(df)}

\CommentTok{\# Create scatter plot}

\NormalTok{plt.figure(figsize}\OperatorTok{=}\NormalTok{(}\DecValTok{9}\NormalTok{, }\FloatTok{5.5}\NormalTok{))}

\NormalTok{plt.scatter(df[}\StringTok{"Study\_Hours"}\NormalTok{], df[}\StringTok{"Actual\_Score"}\NormalTok{], s}\OperatorTok{=}\DecValTok{80}\NormalTok{)}

\NormalTok{plt.xlabel(}\StringTok{"Study Hours"}\NormalTok{, fontsize}\OperatorTok{=}\DecValTok{12}\NormalTok{)}

\NormalTok{plt.ylabel(}\StringTok{"Actual Score"}\NormalTok{, fontsize}\OperatorTok{=}\DecValTok{12}\NormalTok{)}

\NormalTok{plt.title(}\StringTok{"Relationship Between Study Hours and Exam Scores"}\NormalTok{, fontsize}\OperatorTok{=}\DecValTok{14}\NormalTok{, pad}\OperatorTok{=}\DecValTok{15}\NormalTok{)}

\NormalTok{plt.grid(alpha}\OperatorTok{=}\FloatTok{0.2}\NormalTok{)}

\NormalTok{plt.tight\_layout()}

\NormalTok{plt.show()}
\end{Highlighting}
\end{Shaded}

\begin{verbatim}
   Student  Study_Hours  Actual_Score
0        1          1.5            43
1        2          2.0            47
2        3          3.5            56
3        4          4.0            59
4        5          5.5            68
5        6          6.0            70
6        7          7.5            78
7        8          8.0            81
8        9          9.5            87
9       10         10.0            89
\end{verbatim}

\pandocbounded{\includegraphics[keepaspectratio]{2_classical_forecast_tech_pdf_files/figure-pdf/cell-2-output-2.pdf}}

\textbf{:::}

\subsection{What do we see?}\label{what-do-we-see}

The scatter plot suggests a ****positive relationship**** between study
hours and examination scores.

Students who study more tend to obtain higher scores.

However, the relationship is ****not perfectly deterministic****.

For example, students who study a similar number of hours do not
necessarily obtain exactly the same score.

This suggests that something systematic is happening, but there is also
some randomness in the data.

\subsection{Signal and Noise}\label{signal-and-noise}

We can think of the observed score as consisting of two components:

\$\$

\boxed{\text{Observed Data} = \text{Signal} + \text{Noise}}

\$\$

*****Signal*****

The ****signal**** is the systematic pattern in the data.

In our example, the signal is the tendency for examination scores to
increase as study hours increase.

*****Noise*****

The ****noise**** is the random variation around the systematic pattern.

Even if two students study for approximately the same amount of time,
their scores may differ because of:

\begin{itemize}
\item
  Prior knowledge
\item
  Individual ability
\item
  Concentration
\item
  Exam difficulty
\item
  Luck
\item
  Other factors
\end{itemize}

Therefore, we do not expect every observation to lie exactly on the
underlying relationship.

\subsection{Capturing the Signal with Linear
Regression}\label{capturing-the-signal-with-linear-regression}

One simple way to capture the systematic component of the relationship
is ****linear regression****.

We specify:

\$\$

Score\_i = \beta\_0 + \beta\_1 StudyHours\_i + \varepsilon\_i

\$\$

where:

\begin{itemize}
\item
  \(\beta_0\) represents the intercept
\item
  \(\beta_1\) represents the effect of study hours
\item
  \(\varepsilon_i\) represents the noise
\end{itemize}

The regression model attempts to identify the ****signal**** while
treating the remaining variation as noise.

In estimated form:

\$\$

\widehat{Score}\_i

=

\hat{\beta}\_0

+

\hat{\beta}\_1 StudyHours\_i

\$\$




\end{document}
